\section{Running experiments: the sweep driver}
\label{sec:platform-sweep}

The harness makes one \texttt{(method, problem, seed)} run work. An experiment
is many such runs over a grid (seeds, problem sizes, training distributions),
with results in one table, and the bookkeeping that entails is the same for
every experiment. It is therefore factored into
\texttt{experiments/sweep.py}, and every experiment in this thesis is a short
script on top of it. Experiments are deliberately kept outside the core: they
are consumers of the platform, not part of its interface, so adding an
experiment never touches the code that other experiments depend on.

An experiment declares a \texttt{SweepSpec}: the axes of the grid
(\texttt{Axis("seed", ...)}, \texttt{Axis("method", ...)}, \ldots), optional
\emph{prepare} hooks that run once per prefix of the grid (for example,
generate the dataset once per seed and reuse it for every method), and the
function \texttt{run\_cell} that turns one grid point into one or more result
rows. The driver supplies everything else:
\begin{itemize}
  \item \emph{Resume.} Cells whose rows already exist in \texttt{results.jsonl}
    are skipped, so a sweep that stops at hour 30 of 40 restarts at hour 30.
  \item \emph{Error isolation.} A cell that raises is written to
    \texttt{errors.jsonl} with its traceback and the sweep continues; a
    \texttt{stop\_on\_error} switch inverts this while developing.
  \item \emph{Correct axis fields.} The driver stamps each row with the axis
    values it actually ran, so a table cannot record a seed the run did not
    use.
  \item \emph{Provenance}, as described in Section~\ref{sec:platform-repro}.
  \item \emph{Summaries.} Every numeric field of every row is aggregated
    across seeds into \texttt{summary.json}; the statistics are a parameter
    (mean and standard deviation by default, extended rather than replaced),
    and a metric added to a row appears in the summary without further code.
  \item \emph{Command-line overrides} derived from the experiment's own
    configuration dataclass, so a tunable is declared once.
\end{itemize}
What the driver deliberately does not own is the experiment's configuration
and the translation of its hyperparameters into a library's flags. Flag names
differ per library, and two experiments holding the same value for different
reasons are not duplication: the reasons are the content, and they stay in
each experiment's file.

For goal (iii) the concrete deliverable is
\texttt{experiments/experiment\_template.py}: a complete, runnable sweep
(LISTA on a small LASSO dataset, a couple of minutes on a CPU) with four
marked places to edit, namely the configuration dataclass, the prepare hook,
the flag translation and \texttt{run\_cell}. A newcomer copies the template,
edits those four places, and inherits resume, error isolation, provenance
and summaries without writing any of them. The accompanying
\texttt{CONTRIBUTING.md} walks through adding a new method the same way,
assuming no prior knowledge of the repository.
